\paragraph{Decision objectives and grader errors.}
A separate 500-repetition ablation on each of six scenarios compares
success-only, component-only, guarded-win, and fixed weighted-utility
rules. These test different objectives. With identical agents,
component-only noninferiority qualifies 78.4\% of streams while guarded
win qualifies 0.2\%: satisfying component margins does not establish a
preference improvement. The fixed utility $0.8S+0.1C+0.1/(1+\mathrm{cost})$
approves the compliance-regression scenario in 32.6\% of runs, correctly
reflecting its positive utility target; adding component requirements
reduces this to zero. No rule is uniformly more powerful across different
objectives. Guarded bounded-efficiency superiority matches guarded win
in both deployment rate and mean capped pair count in the efficiency-gain
and joint-gain cases; these cases show no decision-speed advantage for
the hierarchical objective over that component-based rule.
Appendix~\ref{app:ablations} explains the objectives; the supplementary
CSV tables retain all comparisons and population Pareto descriptions.

Four additional cases perturb success grading. With 5\% false-positive
errors for $A$ and 5\% false-negative errors for $B$, a true success
contrast of $-0.04$ becomes $+0.012$ measured. Measured-label guarded
inference approves 489/500 streams (97.8\%), while true-label inference
approves none. A conservative rule using the specified error bounds also
approves none. Conversely, adverse differential grading prevents detection
of a genuine efficiency gain. Optional-stopping validity for measured
outcomes cannot repair a biased definition of success. The error bounds
are known simulation inputs, not empirically validated grader guarantees.
